\documentclass[10pt,a4paper]{amsart}
\usepackage{amsmath,amssymb,amsthm,graphicx}
% Retained preamble for the Stacks Project chapter {CH} (shared project preamble).
% The chapter file pinned-original-source.tex only carries \input{{preamble}} on its first line,
% so no preamble text lives inside the chapter itself.  The Stacks Project uses one single shared
% preamble for every chapter: the snapshot stages it three times (sources/stacks-chow/preamble.tex,
% sources/stacks-etale-cohomology/preamble.tex and focused-seeds/AGP/sources/stacks-snapshot-20261011/
% preamble.tex) and those copies are identical in size.  The block below is copied mechanically,
% line for line, from sources/stacks-chow/preamble.tex lines 14-18 and 154-217: the Xy-pic setup and
% the native theorem declarations together with the mathematical macros.  Nothing was added, removed
% or redefined; the wrapper class and the amsmath/amssymb/amsthm/graphicx base packages are supplied
% by the assembler.
% For commutative diagrams we use Xy-pic
% Chapter: stacks-cohomology.tex -- Cohomology of Sheaves.
\usepackage[all]{xy}

% We use 2cell for 2-commutative diagrams.
\xyoption{2cell}

% Theorem environments.
%
\theoremstyle{plain}
\newtheorem{theorem}[subsection]{Theorem}
\newtheorem{proposition}[subsection]{Proposition}
\newtheorem{lemma}[subsection]{Lemma}

\theoremstyle{definition}
\newtheorem{definition}[subsection]{Definition}
\newtheorem{example}[subsection]{Example}
\newtheorem{exercise}[subsection]{Exercise}
\newtheorem{situation}[subsection]{Situation}

\theoremstyle{remark}
\newtheorem{remark}[subsection]{Remark}
\newtheorem{remarks}[subsection]{Remarks}

\numberwithin{equation}{subsection}

% Macros
%
\def\lim{\mathop{\mathrm{lim}}\nolimits}
\def\colim{\mathop{\mathrm{colim}}\nolimits}
\def\Spec{\mathop{\mathrm{Spec}}}
\def\Hom{\mathop{\mathrm{Hom}}\nolimits}
\def\Ext{\mathop{\mathrm{Ext}}\nolimits}
\def\SheafHom{\mathop{\mathcal{H}\!\mathit{om}}\nolimits}
\def\SheafExt{\mathop{\mathcal{E}\!\mathit{xt}}\nolimits}
\def\Sch{\mathit{Sch}}
\def\Mor{\mathop{\mathrm{Mor}}\nolimits}
\def\Ob{\mathop{\mathrm{Ob}}\nolimits}
\def\Sh{\mathop{\mathit{Sh}}\nolimits}
\def\NL{\mathop{N\!L}\nolimits}
\def\CH{\mathop{\mathrm{CH}}\nolimits}
\def\proetale{{pro\text{-}\acute{e}tale}}
\def\etale{{\acute{e}tale}}
\def\QCoh{\mathit{QCoh}}
\def\Ker{\mathop{\mathrm{Ker}}}
\def\Im{\mathop{\mathrm{Im}}}
\def\Coker{\mathop{\mathrm{Coker}}}
\def\Coim{\mathop{\mathrm{Coim}}}

% Boxtimes
%
\DeclareMathSymbol{\boxtimes}{\mathbin}{AMSa}{"02}

%
% Macros for moduli stacks/spaces
%
\def\QCohstack{\mathcal{QC}\!\mathit{oh}}
\def\Cohstack{\mathcal{C}\!\mathit{oh}}
\def\Spacesstack{\mathcal{S}\!\mathit{paces}}
\def\Quotfunctor{\mathrm{Quot}}
\def\Hilbfunctor{\mathrm{Hilb}}
\def\Curvesstack{\mathcal{C}\!\mathit{urves}}
\def\Polarizedstack{\mathcal{P}\!\mathit{olarized}}
\def\Complexesstack{\mathcal{C}\!\mathit{omplexes}}
% \Pic is the operator that assigns to X its picard group, usage \Pic(X)
% \Picardstack_{X/B} denotes the Picard stack of X over B
% \Picardfunctor_{X/B} denotes the Picard functor of X over B
\def\Pic{\mathop{\mathrm{Pic}}\nolimits}
\def\Picardstack{\mathcal{P}\!\mathit{ic}}
\def\Picardfunctor{\mathrm{Pic}}
\def\Deformationcategory{\mathcal{D}\!\mathit{ef}}

% --- appended 2026-10-12 (Chief): the project own definitions of the reference/slogan/history
% display environments, copied mechanically from the shared Stacks preamble lines 8-12
% (\usepackage{verbatim} plus the three \newenvironment lines, which the chapter preamble keeps
% verbatim from the project preamble). Nothing added, removed or redefined: these environments
% carry no mathematical content and the project itself typesets them as comments. Without these
% three lines a counted statement whose own text uses them fails with Environment slogan undefined.
\usepackage{verbatim}
\newenvironment{reference}{\comment}{\endcomment}
\newenvironment{slogan}{\comment}{\endcomment}
\newenvironment{history}{\comment}{\endcomment}
\usepackage{cleveref}
\crefname{originalequation}{original equation}{original equations}
\Crefname{originalequation}{Original equation}{Original equations}
\makeatletter
\newlabel{equation-global-sections-to-alternating-cech}{{08C0}{source}{}{sourcepointer-1}{}}
\newlabel{equation-global-sections-to-alternating-cech@cref}{{[originalequation][0][]08C0}{[1][0][]source}{}{}{}}
\makeatother
\crefname{originalreference}{original reference}{original references}
\Crefname{originalreference}{Original reference}{Original references}
\makeatletter
\newlabel{homology-lemma-double-complex-gives-resolution}{{0133}{source}{}{sourcepointer-2}{}}
\newlabel{homology-lemma-double-complex-gives-resolution@cref}{{[originalreference][0][]0133}{[1][0][]source}{}{}{}}
\makeatother
\crefname{originallemma}{original lemma}{original lemmas}
\Crefname{originallemma}{Original lemma}{Original lemmas}
\makeatletter
\newlabel{lemma-injective-trivial-cech}{{01EP}{source}{}{sourcepointer-3}{}}
\newlabel{lemma-injective-trivial-cech@cref}{{[originallemma][0][]01EP}{[1][0][]source}{}{}{}}
\makeatother
\crefname{originallemma}{original lemma}{original lemmas}
\Crefname{originallemma}{Original lemma}{Original lemmas}
\makeatletter
\newlabel{lemma-alternating-usual}{{01FM}{source}{}{sourcepointer-4}{}}
\newlabel{lemma-alternating-usual@cref}{{[originallemma][0][]01FM}{[1][0][]source}{}{}{}}
\makeatother
\crefname{originalreference}{original reference}{original references}
\Crefname{originalreference}{Original reference}{Original references}
\makeatletter
\newlabel{injectives-theorem-K-injective-embedding-grothendieck}{{079P}{source}{}{sourcepointer-5}{}}
\newlabel{injectives-theorem-K-injective-embedding-grothendieck@cref}{{[originalreference][0][]079P}{[1][0][]source}{}{}{}}
\makeatother
\title{Construction of the canonical functorial map from the alternating Cech complex of a finite covering to derived global sections}
\author{The Stacks Project Authors}
\date{stacks-project:cohomology.tex:local-snapshot-20261011; GNU FDL 1.2 or later}
\begin{document}
\maketitle
\noindent\textit{Editorial note.} All delivered mathematical text is the author's own, copied verbatim from the pinned Stacks Project chapter file cohomology.tex; no statement or proof step was rewritten, reordered or removed.\par\medskip
\noindent\textit{Editorial note.} The counted span cites six cross-chapter labels (including a theorem), each recorded in assembly.external\_references with its canonical tag as a navigation pointer.\par\medskip
\noindent\textit{Editorial note.} The author's closing sentence omits only the verification of functoriality and of the compatibilities of the constructed map; that sentence is kept verbatim in the proof and the rest of the argument is complete.\par\medskip
\noindent\textit{Allowed standard prerequisites (editorial).}\begin{itemize}
\item Existence of a K-injective embedding. Stacks Project, Injectives, Theorem injectives-theorem-K-injective-embedding-grothendieck (tag 079P). Supplies the K-injective representative of F\textasciicircum{}bullet used to define the map.
\item Comparison of the alternating and usual Cech complexes and injective triviality. Stacks Project, Cohomology of Sheaves, Lemmas lemma-injective-trivial-cech (tag 08C2) and lemma-alternating-usual (tag 01FM). Provide the comparison that turns the levelwise map into the required map into derived global sections.
\end{itemize}\medskip
\ifdefined\hypertarget\hypertarget{sourcepointer-1}{}\fi
\noindent\textit{Source reference.} Original equation 08C0: sources/stacks-chow/tags.txt, source lines 10585--10585. Cross-chapter standard prerequisite cited by the author inside the counted span; it is used as a pointer and its own proof is not part of this item.\par\medskip
\ifdefined\hypertarget\hypertarget{sourcepointer-2}{}\fi
\noindent\textit{Source reference.} Original reference 0133: sources/stacks-chow/tags.txt, source lines 1421--1421. Cross-chapter standard prerequisite cited by the author inside the counted span; it is used as a pointer and its own proof is not part of this item.\par\medskip
\ifdefined\hypertarget\hypertarget{sourcepointer-3}{}\fi
\noindent\textit{Source reference.} Original lemma 01EP: sources/stacks-chow/tags.txt, source lines 1842--1842. Cross-chapter standard prerequisite cited by the author inside the counted span; it is used as a pointer and its own proof is not part of this item.\par\medskip
\ifdefined\hypertarget\hypertarget{sourcepointer-4}{}\fi
\noindent\textit{Source reference.} Original lemma 01FM: sources/stacks-chow/tags.txt, source lines 1875--1875. Cross-chapter standard prerequisite cited by the author inside the counted span; it is used as a pointer and its own proof is not part of this item.\par\medskip
\ifdefined\hypertarget\hypertarget{sourcepointer-5}{}\fi
\noindent\textit{Source reference.} Original reference 079P: sources/stacks-chow/tags.txt, source lines 9279--9279. Cross-chapter standard prerequisite cited by the author inside the counted span; it is used as a pointer and its own proof is not part of this item.\par\medskip
\begin{lemma}
\label{lemma-alternating-cech-complex-complex}
Let $(X, \mathcal{O}_X)$ be a ringed space.
Let $\mathcal{U} : X = \bigcup_{i \in I} U_i$ be
a finite open covering. For a complex $\mathcal{F}^\bullet$
of $\mathcal{O}_X$-modules there is a canonical map
$$
\text{Tot}(\check{\mathcal{C}}^\bullet_{alt}(\mathcal{U}, \mathcal{F}^\bullet))
\longrightarrow
R\Gamma(X, \mathcal{F}^\bullet)
$$
functorial in $\mathcal{F}^\bullet$ and compatible with
(\ref{equation-global-sections-to-alternating-cech}).
\end{lemma}

\begin{proof}
Let ${\mathcal I}^\bullet$ be a K-injective complex whose terms
are injective $\mathcal{O}_X$-modules.
The map (\ref{equation-global-sections-to-alternating-cech}) for
$\mathcal{I}^\bullet$ is a map
$\Gamma(X, {\mathcal I}^\bullet) \to
\text{Tot}(\check{\mathcal{C}}^\bullet_{alt}({\mathcal U},
{\mathcal I}^\bullet))$.
This is a quasi-isomorphism of complexes of abelian groups
as follows from
Homology, Lemma \ref{homology-lemma-double-complex-gives-resolution}
applied to the double complex
$\check{\mathcal{C}}^\bullet_{alt}({\mathcal U},
{\mathcal I}^\bullet)$ using
Lemmas \ref{lemma-injective-trivial-cech} and \ref{lemma-alternating-usual}.
Suppose ${\mathcal F}^\bullet \to {\mathcal I}^\bullet$ is a quasi-isomorphism
of ${\mathcal F}^\bullet$ into a K-injective complex whose terms
are injectives (Injectives, Theorem
\ref{injectives-theorem-K-injective-embedding-grothendieck}).
Since $R\Gamma(X, {\mathcal F}^\bullet)$ is represented by the complex
$\Gamma(X, {\mathcal I}^\bullet)$ we obtain the map of the lemma
using
$$
\text{Tot}(\check{\mathcal{C}}^\bullet_{alt}({\mathcal U},
{\mathcal F}^\bullet))
\longrightarrow
\text{Tot}(\check{\mathcal{C}}^\bullet_{alt}({\mathcal U},
{\mathcal I}^\bullet)).
$$
We omit the verification of functoriality and compatibilities.
\end{proof}

\end{document}
